\chapter{Introductie}
\label{chap:introductie}

\begin{wrapfigure}{r}{0.5\textwidth}[10]
  \centering
  \includegraphics[width=0.6\textwidth]{image1.png}
  \caption{}
  \label{fig:image1.png}
 \end{wrapfigure}

We leggen deze figuur uit met een eenvoudig voorbeeld.
Je hebt een doodgewone knop die je kan indrukken.
Jij als \textbf{operator} geef controle inputs aan de knop, 
bijvoorbeeld door hem in te drukken.

\textbf{De logic controller (PLC) }ontvangt deze input en 
verwerkt deze volgens een vooraf gedefinieerd programma.

De knop krijgt info door sensors die bijvoorbeeld info geven of 
de knop al dan niet is ingedrukt van de machine

Als de knop is ingedrukt gaat door er een \textbf{logic actions} gebeuren waarbij de \textbf{machine} een actie uitvoert,
bijvoorbeeld een motor starten of een lamp laten branden.
en de \textbf{operator} krijgt feedback van de \textbf{machine}, bijvoorbeeld door een lamp die gaat branden of een geluidssignaal.

